\documentclass[11pt,a4paper]{article}

\usepackage[utf8]{inputenc}
\usepackage[T1]{fontenc}
\usepackage{lmodern}
\usepackage[margin=0.78in]{geometry}
\usepackage[table]{xcolor}
\usepackage{array}
\usepackage{tabularx}
\usepackage{booktabs}
\usepackage{enumitem}
\usepackage{titlesec}
\usepackage{fancyhdr}
\usepackage{lastpage}
\usepackage{hyperref}

\definecolor{TextInk}{HTML}{233142}
\definecolor{HeadingBlue}{HTML}{1D3557}
\definecolor{AccentTeal}{HTML}{2A7F92}
\definecolor{AccentGold}{HTML}{C58B4E}
\definecolor{TableStripe}{HTML}{F3F8FA}
\definecolor{SoftBlue}{HTML}{EEF6F8}
\definecolor{SoftGreen}{HTML}{EEF7F0}
\definecolor{SoftGold}{HTML}{FCF6E8}

\hypersetup{
  colorlinks=true,
  linkcolor=HeadingBlue,
  urlcolor=AccentTeal,
  pdftitle={May 2 Paper 1 Simulation Teacher Marking Notes},
  pdfauthor={IB Geography SL Preparation}
}

\color{TextInk}
\setlength{\parskip}{0.45em}
\setlength{\parindent}{0pt}
\setlength{\tabcolsep}{5pt}
\renewcommand{\arraystretch}{1.18}
\setlist[itemize]{leftmargin=1.35em,itemsep=3pt,topsep=4pt}

\newcolumntype{Y}{>{\raggedright\arraybackslash}X}
\newcolumntype{L}[1]{>{\raggedright\arraybackslash}p{#1}}
\newcommand{\term}[1]{\textcolor{HeadingBlue}{\textbf{#1}}}
\newcommand{\score}[1]{\textcolor{AccentTeal}{\textbf{#1}}}
\newcommand{\boxnote}[3]{%
\noindent\colorbox{#1}{%
  \begin{minipage}{0.965\textwidth}
  \sffamily
  \textbf{\textcolor{HeadingBlue}{#2}} #3
  \end{minipage}
}}

\titleformat{\section}
  {\normalfont\Large\sffamily\bfseries\color{HeadingBlue}}
  {}{0pt}{}
  [\vspace{0.08em}{\color{AccentGold}\titlerule[0.8pt]}]

\titleformat{\subsection}
  {\normalfont\large\sffamily\bfseries\color{AccentTeal}}
  {}{0pt}{}

\pagestyle{fancy}
\setlength{\headheight}{14pt}
\fancyhf{}
\fancyhead[L]{\sffamily\color{AccentTeal}Teacher Marking Notes}
\fancyhead[R]{\sffamily\color{HeadingBlue}May 2 Paper 1 Simulation}
\fancyfoot[C]{\sffamily\color{HeadingBlue}\thepage\ of \pageref{LastPage}}
\renewcommand{\headrulewidth}{0.4pt}

\begin{document}

\begin{center}
  {\sffamily\bfseries\color{HeadingBlue}\LARGE May 2 Paper 1 Simulation Teacher Marking Notes}\\[0.25em]
  {\sffamily\color{AccentTeal}\large Urban Environments Option}
\end{center}

\boxnote{SoftBlue}{Basis for marking.}{This marking uses the May 2 marking
guide and the timed student transcript in the \texttt{plans} folder. The
student answered Question 4(b).}

\section{Overall Mark}

\rowcolors{2}{TableStripe}{white}
\begin{tabularx}{\textwidth}{L{0.18\textwidth}L{0.14\textwidth}Y}
\toprule
\textbf{Question} & \textbf{Mark} & \textbf{Teacher rationale} \\
\midrule
Q1 & \score{1/2} & Identifies uneven green space, which is relevant to
environmental quality. The second point, uneven population density, is not
clearly environmental quality evidence, and no source data or zones are named.
This is a lenient award; a stricter source-evidence reading could score
\score{0/2}. \\
Q2 & \score{4/4} & Two clear contrasts are given: rent index
(\term{82} versus \term{24}) and green space (\term{14\%} versus \term{3\%}).
Both contrasts identify the two zones and use source evidence. \\
Q3 & \score{4/4} & Two valid urban-stress chains are explained: urban sprawl
leading to longer commuting, congestion, and pollution; and rapid growth
outpacing infrastructure, increasing informal settlement growth. \\
Q4(b) & \score{7/10} & Reaches the 7--8 band. The answer is balanced, uses
Detroit as a named example, considers benefits and costs, and includes a clear
conclusion. It remains at the lower end of the band because the Detroit evidence
is sometimes general, the regeneration process is not supported with precise
projects or data, and the Rio comparison is useful but not fully integrated into
the main Detroit argument. \\
\midrule
\textbf{Total} & \score{16/20} & Strong performance. Main lost marks come from
Q1 precision and Q4 case-study depth. \\
\bottomrule
\end{tabularx}
\rowcolors{2}{}{}

\section{Question-By-Question Feedback}

\subsection{Q1: State Evidence Of Uneven Environmental Quality}

Award \score{1/2} leniently. The answer gives one valid idea: green space is uneven.
However, a full-mark answer needed two specific pieces of source evidence, for
example \term{34\%} green space in the suburban zone compared with \term{3\%}
in the peripheral informal settlement. Population density may affect urban
pressure, but it is not by itself evidence of environmental quality unless
linked to a clear environmental condition. Under a stricter application of the
``source-based point'' rule, this response could receive \score{0/2}.

\subsection{Q2: Describe Two Contrasts}

Award \score{4/4}. This is the strongest source-use answer. The student
correctly compares inner-city regeneration zones with peripheral informal
settlements using rent and green-space data. The answer satisfies the command
term because it describes contrasts rather than drifting into explanation.

\subsection{Q3: Explain Urban Stress}

Award \score{4/4}. The response gives two cause-and-effect chains, which matches
the marking guide. The grammar is imperfect, but the geography is clear: rapid
growth can increase commuting pressure, congestion, pollution, infrastructure
stress, and informal settlement growth.

\subsection{Q4(b): Regeneration Or Gentrification}

Award \score{7/10}. This is a clear evaluative answer. It explains that
regeneration and gentrification can bring investment, services, jobs, business
activity, and renewed attractiveness, but can also raise living costs, displace
residents, and create racial or social tensions.

The Detroit case study is relevant and has useful historical context:
automobile growth, population decline, suburbanisation, abandoned housing, crime,
and reinvestment in downtown and Midtown. The answer also makes a valid wider
comparison to Rio de Janeiro and favela displacement before the 2016 Olympics.

To move into \score{8--10}, the answer needs more precise case-study evidence.
For example, it could name specific Detroit regeneration projects, give clearer
before-and-after evidence, and judge the scale of success more explicitly. The
Rio example should either be integrated as a comparison or omitted to keep the
argument tightly focused.

\section{Summary}

\boxnote{SoftGreen}{Strengths.}{The student understands the main urban
processes and writes in causal chains. Q2 and Q3 are secure. The 10-mark answer
has balance, a named example, and a conclusion, which places it in a strong
band.}

\boxnote{SoftGold}{Priority improvement.}{The student should focus on
precision: use exact source evidence in short questions and exact case-study
evidence in 10-mark answers. The ideas are present, but some marks are lost
because the evidence is too general.}

\section{Recommendations}

\begin{itemize}
  \item For \term{state} questions, write two compact source facts with numbers
  and zones. A reliable structure is: ``Zone A has X, whereas Zone B has Y.''
  \item For \term{describe} questions, keep doing what worked in Q2: name both
  places or categories, then include the data.
  \item For \term{explain} questions, keep using cause-and-effect chains. The
  Q3 structure is effective.
  \item For 10-mark answers, prepare three Detroit evidence anchors: one for
  decline, one for regeneration, and one for uneven impacts or displacement.
  \item Add a more explicit evaluative sentence at the end of each body
  paragraph, such as: ``This solves the problem for investors and commuters, but
  not for low-income residents because...''
  \item In the final conclusion, answer the command term directly with degree:
  regeneration solves some visible symptoms of urban decline, but only partly
  solves deeper inequality because benefits are spatially and socially uneven.
\end{itemize}

\section{Next Practice Task}

Rewrite only Q1 and Q4(b).

\begin{itemize}
  \item Q1 target: \score{2/2} using two source statistics.
  \item Q4(b) target: \score{8/10} by adding named Detroit project evidence,
  clearer scale, and one sharper judgment about who benefits and who loses.
\end{itemize}

\end{document}
